../cictikz/src/cictikz/data/tex/cictikz_preamble.tex

% The vibration modes of a quartz blank, redrawn to answer the question
% the picture is really there for: which mode gives which frequency.
%
% Each panel shows the blank at rest as a solid outline and one extreme
% of its motion dashed on top, with red arrows for the direction the
% material moves. Underneath is the range that mode is useful over,
% because the mode is not a curiosity - it is what decides whether a
% crystal is a 32 kHz timekeeper or a 40 MHz radio reference.

\begin{tikzpicture}[thick]

  \def\pw{4.0}       % panel pitch
  \def\yc{2.0}       % where a blank sits inside its panel
  \def\ylabel{0.2}   % mode name
  \def\yuse{-0.35}   % what it is good for

  \tikzset{
    rest/.style   = {draw=gray!70, fill=oxteal!30},
    moved/.style  = {draw=black, dashed, thick},
    motion/.style = {red, line width=1.1pt, -{Latex[length=2.6mm]}},
  }

  % panel #1: index, #2: mode name, #3: what it is used for
  \newcommand{\modelabel}[3]{
    \node[anchor=north, scale=0.95] at ({#1*\pw+1.6},\ylabel) {#2};
    \node[anchor=north, scale=0.8, gray!50!black, align=center]
      at ({#1*\pw+1.6},\yuse) {#3};
  }

  % ================= 1. longitudinal =================
  \begin{scope}[shift={(0,0)}]
    \fill[rest] (0.35,{\yc-0.6}) rectangle (2.85,{\yc+0.6});
    \draw[moved] (0.05,{\yc-0.45}) rectangle (3.15,{\yc+0.45});
    \draw[motion] (1.6,\yc) -- (0.35,\yc);
    \draw[motion] (1.6,\yc) -- (2.85,\yc);
  \end{scope}
  \modelabel{0}{Longitudinal}{the bar stretches\\ and shortens\\ tens of kHz}

  % ================= 2. thickness shear =================
  \begin{scope}[shift={(\pw,0)}]
    \fill[rest] (0.35,{\yc-0.6}) rectangle (2.85,{\yc+0.6});
    \draw[moved] (0.75,{\yc+0.6}) -- (3.25,{\yc+0.6})
      -- (2.45,{\yc-0.6}) -- (-0.05,{\yc-0.6}) -- cycle;
    \draw[motion] (1.2,{\yc+0.3}) -- (2.4,{\yc+0.3});
    \draw[motion] (2.0,{\yc-0.3}) -- (0.8,{\yc-0.3});
  \end{scope}
  \modelabel{1}{Thickness shear}{faces slide past\\ each other\\ 1 to 200 MHz}

  % ================= 3. flexural =================
  \begin{scope}[shift={(2*\pw,0)}]
    \fill[rest] (0.35,{\yc-0.5}) rectangle (2.85,{\yc+0.5});
    \draw[moved] (0.35,{\yc-0.5})
      .. controls (1.1,{\yc+0.25}) and (2.1,{\yc+0.25}) .. (2.85,{\yc-0.5})
      -- (2.85,{\yc+0.5})
      .. controls (2.1,{\yc+1.25}) and (1.1,{\yc+1.25}) .. (0.35,{\yc+0.5})
      -- cycle;
    \draw[motion] (1.6,{\yc+0.55}) -- (1.6,{\yc+1.25});
    \draw[motion] (0.6,{\yc+0.4}) -- (0.6,{\yc-0.3});
    \draw[motion] (2.6,{\yc+0.4}) -- (2.6,{\yc-0.3});
  \end{scope}
  \modelabel{2}{Flexural}{the bar bends\\ like a diving board\\ tens of kHz}

  % ================= 4. face shear =================
  \begin{scope}[shift={(3*\pw,0)}]
    \fill[rest] (0.65,{\yc-0.85}) rectangle (2.55,{\yc+0.85});
    \draw[moved] (0.95,{\yc+0.85}) -- (2.85,{\yc+0.55})
      -- (2.25,{\yc-0.85}) -- (0.35,{\yc-0.55}) -- cycle;
    \draw[motion] (1.15,{\yc+0.5}) -- (2.15,{\yc+0.5});
    \draw[motion] (2.05,{\yc-0.5}) -- (1.05,{\yc-0.5});
  \end{scope}
  \modelabel{3}{Face shear}{the square racks\\ over into a rhombus\\ hundreds of kHz}

  % ================= 5. tuning fork =================
  \begin{scope}[shift={(4*\pw,0)}]
    \fill[rest] (0.25,{\yc-0.55}) rectangle (1.15,{\yc+0.55});
    \fill[rest] (1.15,{\yc+0.12}) rectangle (2.85,{\yc+0.45});
    \fill[rest] (1.15,{\yc-0.45}) rectangle (2.85,{\yc-0.12});
    \draw[moved] (1.15,{\yc+0.12}) -- (2.9,{\yc+0.55})
      -- (2.9,{\yc+0.88}) -- (1.15,{\yc+0.45}) -- cycle;
    \draw[moved] (1.15,{\yc-0.12}) -- (2.9,{\yc-0.55})
      -- (2.9,{\yc-0.88}) -- (1.15,{\yc-0.45}) -- cycle;
    \draw[motion] (3.2,{\yc+0.1}) -- (3.2,{\yc+0.85});
    \draw[motion] (3.2,{\yc-0.1}) -- (3.2,{\yc-0.85});
  \end{scope}
  \modelabel{4}{Tuning fork}{two tines swing\\ apart and together\\ 32.768 kHz}

  % ---- the two modes an integrated circuit actually meets ----
  \foreach \k in {1,4} {
    \draw[armygreen, thin, rounded corners=3pt]
      ({\k*\pw-0.25},-1.55) rectangle ({\k*\pw+3.45},{\yc+1.35});
  }
  \node[armygreen, anchor=north, scale=0.85, align=center] at ({2*\pw+1.6},-1.8)
    {The two in the green boxes are the ones you will meet on a circuit board:
     a thickness shear blank for the megahertz reference,\\
     and a tuning fork for the 32.768 kHz real-time clock};

  % ---- the point of the picture ----
  \node[anchor=south, align=center, scale=0.95] at ({2*\pw+1.6},{\yc+1.7})
    {A quartz blank can vibrate in any of these. The cut of the blank
     and the circuit around it pick one,\\
     and that choice, not the quartz, is what sets the frequency};

\end{tikzpicture}
\end{document}
